% Section 1.3, from lectures/L01/README.md: what you should be able to do, the questions to test
% yourself with, and the next lecture.

\section{Sammanfattning}
\label{c1:sec:summary}

\needspace{6\baselineskip}
\subsection*{Det här ska du kunna nu}
Efter det här kapitlet ska du kunna:
\begin{itemize}
  \item Förklara skillnaden mellan en digital och en analog signal, och varför digitala signaler
    tål brus.
  \item Beskriva det decimala, det binära och det hexadecimala talsystemet: talbas, siffror och
    positionernas vikter.
  \item Omvandla tal mellan decimalt, binärt och hexadecimalt för hand, i alla riktningar.
  \item Ange hur många olika värden ett tal med $n$ bitar kan anta, och vilket det största är.
  \item Addera två binära tal, och avgöra om summan ryms i åtta bitar.
  \item Beskriva BCD, Graykod, ASCII och 7-segmentkod, och koda och avkoda tal och tecken i dem.
  \item Beräkna en jämn paritetsbit, och förklara vilka fel den upptäcker.
\end{itemize}

\needspace{6\baselineskip}
\subsection*{Frågor att testa dig själv med}
\begin{enumerate}
  \item Varför har den minst signifikanta positionen vikten 1 i \emph{alla} talsystem?
  \item Hur ser du direkt på ett binärt tal om det är jämnt eller udda?
  \item Varför används hexadecimala tal till att skriva binära tal, och inte decimala?
  \item Talet \code{0x100} är ett större tal än \code{0xFF}, men bara med ett. Hur många bitar
    behövs för var och ett av dem?
  \item Vad händer med ett binärt tal när du lägger till en nolla längst till höger? Jämför med
    vad som händer med ett decimalt tal.
  \item Tecknet \code{'7'} och talet 7 är olika saker i en dator. Vad skiljer dem åt, och hur gör
    du om det ena till det andra?
  \item I vilken kod ändras exakt en bit mellan två på varandra följande tal, och varför är det
    användbart?
\end{enumerate}

Talsystemen används i resten av boken, och de sitter först när du har räknat många omvandlingar
själv. Räkna hellre för många än för få. Kalkylatorn i Windows, i läget \textbf{Programmerare},
visar samma tal decimalt, binärt och hexadecimalt samtidigt. Använd den för att kontrollera dina
svar, aldrig för att ta fram dem.

\needspace{6\baselineskip}
\subsection*{Nästa kapitel}
Kapitel~\ref{c2:ch} handlar om logiska grindar och boolesk algebra:
\begin{itemize}
  \item Logiska grindar: NOT, AND, OR, NAND, NOR, XOR och XNOR, och deras sanningstabeller.
  \item Boolesk algebra: notationen, räknelagarna och De Morgans lagar.
  \item Från sanningstabell till grindnät, och från grindnät tillbaka till sanningstabell.
  \item Samma logik byggd med strömkontakter på 24~V, inför Labb~1.
\end{itemize}
